\section{جمع‌بندی: انتخاب بهترین روش کاهش رزولوشن و درونیابی}

در این بخش، با توجه به نتایج به‌دست‌آمده از مراحل پیشین، به ارزیابی کلی الگوریتم‌ها پرداخته و بهترین ترکیب برای انجام عملیات کاهش ابعاد (\lr{Downsampling}) و بازسازی (\lr{Upsampling}) را معرفی می‌کنیم.

\subsection*{ارزیابی روش‌های پیش‌پردازش (فیلترهای \lr{Anti-aliasing})}
در فرآیند کاهش ابعاد، ما سه رویکرد مختلف را روی تصاویر اعمال و بررسی کردیم:
\begin{itemize}
    \item \textbf{بدون فیلتر (بخش ۲):} این روش به دلیل عدم رعایت شرط نایکوئیست، منجر به بروز شدیدترین حالت الیاسینگ (ایجاد \lr{Moiré Pattern}) شد و ساختار تصویر به خصوص در نواحی دارای فرکانس‌های بالا (مثل تصویر \lr{Grid}) کاملاً مخدوش گردید.
    \item \textbf{فیلتر ۹ درایه‌ای پیشرفته (بخش ۴):} اگرچه این فیلتر به دلیل عرض زیادِ پنجره توانست الیاسینگ را به صورت کامل مهار کند، اما بخش عمده‌ای از سیگنال اصلی تصویر نیز از بین رفت و خروجی به شدت تار (\lr{Over-blurred}) شد؛ تا جایی که در نرخ ۱۶ ساختار تصویر کاملاً محو گردید و سیگنال مفید از دست رفت.
    \item \textbf{فیلتر ۳ درایه‌ای استاندارد (بخش ۳):} این فیلتر (آرایه \lr{[0.25, 0.5, 0.25]}) بهترین تعادل (\lr{Trade-off}) را میان مهار الیاسینگ و حفظ وضوح تصویر برقرار کرد. فرکانس‌های خطرساز تا حد قابل‌قبولی سرکوب شدند، در حالی که لبه‌های اصلی تصویر همچنان ساختار خود را حفظ کردند.
\end{itemize}

\subsection*{ارزیابی روش‌های درونیابی (\lr{Interpolation})}
برای بازگرداندن تصویر به ابعاد اولیه نیز دو روش مورد ارزیابی قرار گرفت:
\begin{itemize}
    \item \textbf{نزدیک‌ترین همسایه (\lr{Nearest Neighbor}):} این روش لبه‌های تصویر را به شدت دندانه‌دار، بلوکی و پیکسلی کرد و از نظر بصری خروجی نامطلوبی (به ویژه در لبه‌های مورب و منحنی‌ها نظیر چرخ‌وفلک) به همراه داشت.
    \item \textbf{دوخطی (\lr{Bilinear}):} این الگوریتم با میانگین‌گیری فضایی از ۴ پیکسل مجاور، توانست تغییرات شدت روشنایی را نرم‌تر کند و لبه‌های دندانه‌دار را به شکل چشم‌گیری بهبود بخشد.
\end{itemize}

\subsection*{نتیجه‌گیری نهایی}
با استناد به مشاهدات و تحلیل‌های فوق، \textbf{بهترین نتیجه مربوط به ترکیب «کاهش ابعاد با فیلتر ۳ درایه‌ای (مرحله ۳) به همراه درونیابی دوخطی (\lr{Bilinear})» است.} 

دلیل این انتخاب آن است که فیلتر ۳ درایه‌ای به عنوان یک پایین‌گذر ملایم، بدون نابود کردن کامل جزئیات ساختاری، از بروز الگوهای مخرب الیاسینگ جلوگیری می‌کند. در ادامه نیز، درونیابی دوخطی با ایجاد یک گذار نرم میان پیکسل‌ها، از ایجاد نویزهای بلوکی (\lr{Blocking Artifacts}) در تصویر بازسازی‌شده ممانعت به عمل می‌آورد. این ترکیب در هر دو نرخ ۴ و ۱۶، بالاترین کیفیت بصری و بیشترین وفاداری ساختاری به تصویر اصلی را از خود نشان داد.